\documentclass[runningheads]{llncs}
\usepackage[T1]{fontenc}
\usepackage{graphicx}
\begin{document}
\title{MVC Project: Multilayer Perceptron for Student Performance Prediction}
\titlerunning{MLP: Manual to MNIST}
\author{Muhammad Ahmed Butt}
\authorrunning{M.A. Butt}
\institute{National University of Computer and Emerging Sciences (FAST-NU), \\ Islamabad, Pakistan\\
\email{25i3010@nu.edu.pk}\\
\url{https://github.com/MuhammadAhmedButt25i3010/mvc-mlp-25i3010}}

\maketitle 
\begin{abstract}
This report presents a comprehensive implementation and analysis of a Multilayer Perceptron (MLP) starting from mathematical derivation to computer application. This project is divided into two main phases: a manual derivation phase and a programmed execution phase (Python). In the initial phase, a 2-2-2-1 neural architecture was manually constructed to predict student performance. This involved performing step-by-step forward and backward propagation for multiple data samples. Few key insights from the manual are in regards to the vanishing gradient problem being identified in early layers and SGD vs Mini Batch effectiveness.

The second phase transitions these manual concepts into a functional Python implementation. By using the NumPy library, the manual neuron-by-neuron calculations were converted into vectorized matrix operations, allowing the model to scale efficiently. This implementation was specifically tested on the MNIST dataset of handwritten digits. Beyond basic gradient descent, more advanced optimization techniques like Momentum and Adam were integrated to observe their impact on convergence speed and accuracy. The final results are documented through loss curves and performance metrics, demonstrating how a from-scratch MLP can effectively learn complex patterns in high-dimensional image data.

\keywords{Multilayer Perceptron \and Backpropagation \and Gradient Descent \and MNIST \and Neural Networks}
\end{abstract}
\section{Introduction}
\subsection{Multilayer Perceptron (MLP)}
An MLP is a type of feedforward artificial neural network composed of an input layer, one or more hidden layers, and an output layer. Each neuron in a layer is fully connected to neurons in the next layer, with weights determining the strength of these connections and biases allowing flexibility in activation thresholds.

\subsubsection{Project Goal} 
The goal of this project is to use these networks for two different tasks: predicting student success through manual calculations and performing digit classification using Python.
\section{Background}
The MLP architecture used in the manual phase consists of an input layer, two hidden layers, and an output layer (2-2-2-1). Each neuron $j$ in layer $l$ calculates its output as:
\begin{equation}
    a_j^{(l)} = \sigma\left(\sum_{i} w_{ij}^{(l)} a_i^{(l-1)} + b_j^{(l)}\right)
\end{equation}
\begin{equation}
    \sigma(x) = \frac{1}{1 + e^{-x}}
\end{equation}
where $\sigma$ is the Sigmoid activation function. The error is measured using Mean Squared Error (MSE):
\begin{equation}
    L = \frac{1}{N} \sum_{i=1}^{N} (y_i - \hat{y}_i)^2
\end{equation}
\subsection{Vanishing Gradient Problem}
During the manual backpropagation phase, a significant drop in gradient magnitude was observed in the earlier layers ($w_1$–$w_4$). This occurs because the gradient is a product of multiple sigmoid derivatives, which are all $\le 0.25$, causing the update signal to diminish as it propagates toward the input layer.
\section{Experimental Results}
The evaluation was conducted in two parts: a manual trace for conceptual verification and a Python implementation for high-dimensional classification.

\subsection{Manual Trace (Tasks 4-5)}
The manual derivation was performed over three epochs using a shuffled mini-batch schedule with $\eta = 0.1$. The results, summarized in Table~\ref{tab:manual_mse}, confirm that the backpropagation logic effectively reduced the error across all samples despite the vanishing gradient in early layers.

\begin{table}[h]
\centering
\caption{Manual Trace MSE Reduction Schedule}
\label{tab:manual_mse}
\begin{tabular}{|c|c|c|l|c|}
\hline
\textbf{Epoch} & \textbf{Batch} & \textbf{Samples} & \textbf{Action} & \textbf{Recomputed MSE} \\ \hline
1 & 1 & $s(1), s(2)$ & Avg. gradients, update & 0.3022 \\ \hline
1 & 2 & $s(3)$ & Single-sample update & 0.3015 \\ \hline
2 & 1 & $s(3), s(1)$ & Shuffled avg. update & 0.2988 \\ \hline
2 & 2 & $s(2)$ & Single-sample update & 0.2964 \\ \hline
3 & 1 & $s(2), s(3)$ & Shuffled avg. update & 0.2941 \\ \hline
3 & 2 & $s(1)$ & Final single update & \textbf{0.2922} \\ \hline
\end{tabular}
\end{table}

\subsection{Python Implementation and MNIST}
The Python implementation utilized NumPy for vectorization, translating the scalar logic into matrix form. Training on the MNIST dataset revealed that while standard SGD is functional, advanced optimizers like Adam and Momentum significantly accelerated convergence. 

The transition from a 2-2-2-1 network to a larger hidden layer architecture allowed for a high classification accuracy on handwritten digits. The Adam optimizer was particularly effective in overcoming the tiny gradients in the initial layers observed during the manual phase.

\section{Discussion and Conclusion}
The project successfully demonstrated the fundamental mechanics of neural networks. The manual trace provided a deep understanding of weight update dynamics and identified the vanishing gradient problem inherent in sigmoid-based networks. The computational phase validated that vectorization and adaptive learning rates are essential for scaling these mathematical concepts to real-world datasets like MNIST.

\begin{thebibliography}{8}
\bibitem{ref_rumelhart}
Rumelhart, D.E., Hinton, G.E., Williams, R.J.: Learning representations by back-propagating errors. Nature \textbf{323}, 533--536 (1986).

\bibitem{ref_lecun}
LeCun, Y., Bottou, L., Bengio, Y., Haffner, P.: Gradient-based learning applied to document recognition. Proc. of the IEEE \textbf{86}(11), 2278--2324 (1998).

\bibitem{ref_adam}
Kingma, D.P., Ba, J.: Adam: A method for stochastic optimization. arXiv preprint arXiv:1412.6980 (2014).
\end{thebibliography}
\end{document}
